\documentclass[class=NCU_thesis, crop=false]{standalone}
\begin{document}

\chapter{結論與未來展望}

\section{結論}

本研究旨在建立一套由中文文字輸入至虛擬人手語動畫輸出的系統流程，使使用者能透過自然語言輸入取得連續且可觀看之虛擬人手語動畫。
相較於僅提供單詞查詢或播放既有手語影片之系統，本研究進一步整合中文至 gloss sequence 轉換、手語動作資料檢索、EHM-Tracker 追蹤資料處理、動作參數串接與 GUAVA avatar rendering，使中文句子能被轉換為具備人物外觀一致性與動作連續性的虛擬人手語動畫。

在語言層級方面，本研究沿用中文至自然手語翻譯模組，將中文句子轉換為符合台灣自然手語語序之 gloss sequence。
gloss sequence 在本系統中扮演語言層級與動作層級之間的中介表示，使系統能以詞彙為單位查詢對應之手語動作資料。
在動作資料層級方面，本研究將既有手語詞彙影片延伸為可供虛擬人動畫生成使用的動作資料庫，透過 EHM-Tracker 將原始手語影片轉換為包含人體姿態、手部動作與臉部表情等資訊的追蹤參數。
此設計避免完全依賴昂貴的動作捕捉設備，也使既有手語影片資源得以被重新利用於虛擬人動畫生成流程。

在系統設計方面，本研究將手語動作資料建置與使用者動畫生成分為兩類流程，並以固定的追蹤資料目錄作為動作檢索、參數串接與 GUAVA 資料載入之共同介面。
系統透過 gloss 索引將語言輸出對應至詞彙追蹤資料，並在串接前檢查必要資料是否完整；若詞彙缺少可用資料，則回報缺失詞彙並停止後續生成。
此設計使同一筆詞彙動作資料能由不同中文句子與角色重複使用，也使各階段的輸入、輸出及錯誤來源較容易檢查。

在動作生成層級方面，本研究針對多個手語詞彙片段直接串接時容易產生的不連續問題，導入 motion envelope 分析與過渡幀生成方法。
系統會依據手部動作強度估計各詞彙之有效動作區間，降低詞彙前後準備姿勢與收尾姿勢對串接結果的影響。
此外，系統在相鄰詞彙片段之間加入過渡幀，使前一詞彙的結束姿態能逐步銜接至下一詞彙的起始姿態，減少手部位置突變與姿態跳動。
最後，串接後的連續追蹤資料會輸入 GUAVA avatar rendering 模組，產生虛擬人手語動畫輸出。

在詞彙層級參數平滑方面，本研究以 10 個詞彙樣本比較平滑前後之 jerk、normalized jerk cost、spectral arc length 與二維關鍵點對齊結果。
實驗顯示，SMPL-X body pose 與 MANO hand pose normalized jerk cost 平均分別降低 99.90\% 與 99.87\%，且頻譜平滑度結果亦顯示高頻震盪減少，表示此後處理能有效降低追蹤參數中的快速變化。
然而，hand-only PCK@0.10 平均由 0.9476 降至 0.8868，顯示平滑處理雖可改善動作穩定性，但不能解讀為追蹤準確度提升，且手部細節仍需在平滑度與逐幀對齊精度之間取得平衡。

在追蹤資料前處理方面，本研究比較 EHM-Tracker 使用 100-step 與 1000-step 最佳化設定時的差異。
實驗結果顯示，在雙手皆清楚可見且 2D 偵測較穩定的情境下，1000-step 通常能降低 hand reprojection error 並提升 PCK，但其改善幅度整體有限。
相對地，在單手低可見度或手部離開畫面的情境中，較長最佳化步數可能使模型受到低可信度 2D 偵測點影響，導致不可見或低可信度手部被錯誤拉回人體附近，進而造成視覺上的誤追蹤。
因此，本研究在建立後續手語動作資料庫時採用較保守的 100-step 設定，以降低少數嚴重誤追蹤破壞手語語意與動作串接穩定性的風險。

在使用者評估方面，本研究共回收 53 份有效問卷，填答者觀看 25 支系統生成影片後，針對動作自然度、手勢清晰度、整體品質與準確度進行評分。
整體結果顯示，四項指標平均分數介於 3.80 至 4.01 之間，其中準確度與手勢清晰度表現較佳，表示多數填答者認為系統生成動畫大致能呈現原始詞彙影片所表達之內容。
然而，會手語者與不會手語者之評分存在顯著差異，具手語背景者對手形、臉部表情與句子銜接自然度有更高要求，也更容易察覺生成動畫中的細節不足。
此結果顯示，本研究系統已具備中文至虛擬人手語動畫生成之可行性，但若要進一步應用於真實手語溝通或教學情境，仍需強化手部細節、非手部訊息與句子層級自然度。

綜合而言，本研究之主要成果可歸納如下：

\begin{itemize}
    \item \textbf{建立中文至虛擬人手語動畫之模組化系統設計}：
    本研究將手語動作資料建置與使用者動畫生成分為兩類流程，並整合中文至 gloss sequence 轉換、手語動作資料檢索、動作參數串接與 GUAVA 虛擬人渲染，使中文句子能被轉換為連續的虛擬人手語動畫。

    \item \textbf{將既有手語影片資料延伸為可串接之動作資料}：
    本研究透過 EHM-Tracker 將手語詞彙影片轉換為可載入、可查詢與可串接的追蹤資料，建立 gloss 與追蹤資料目錄之索引及完整性檢查方式，降低系統對專業動作捕捉設備的依賴。

    \item \textbf{提出詞彙動作區間擷取、參數平滑與過渡幀串接流程}：
    本研究利用 motion envelope 分析估計手語詞彙之有效動作區間，並透過詞彙層級參數平滑及過渡幀改善多詞彙動畫中的高頻抖動與片段交界不連續問題。量化結果顯示平滑處理能明顯降低 jerk 相關指標，但亦揭示手部平滑度與逐幀對齊精度之間的取捨。

    \item \textbf{分析 EHM-Tracker 最佳化步數對手語動作資料品質之影響}：
    本研究透過 reprojection error、PCK 與視覺檢查比較 100-step 與 1000-step 追蹤結果，指出較長最佳化步數雖可提升部分可見手部的量化指標，但在低可見度手部上可能造成誤追蹤風險。

    \item \textbf{透過使用者問卷驗證系統可行性與不足}：
    本研究以使用者主觀評估分析生成動畫之自然度、清晰度、品質與準確度，結果顯示系統已具備基本可觀看性，同時也揭示手部細節、臉部表情與詞彙銜接仍是後續改善重點。
\end{itemize}

\section{研究限制}

雖然本研究已完成中文至虛擬人手語動畫生成流程，並透過量化實驗與使用者問卷驗證其可行性，但目前系統仍存在若干限制。

首先，本研究之語言層級輸出依賴既有 SLT 模組與手語詞彙資料庫。
當輸入中文句子包含資料庫中不存在的詞彙、複雜語意或需要更細緻語境判斷的表達時，系統可能無法找到對應動作資料，或只能以近似詞彙替代。
因此，目前系統較適合處理已收錄詞彙可覆蓋之句子，對開放式自然語言輸入的泛化能力仍有限。

其次，本研究之手語動作主要來自既有單詞影片。
不同影片之拍攝角度、人物身形、動作速度、起始姿勢與結束姿勢並不完全一致，即使經過 motion envelope 擷取與過渡幀生成，詞彙片段之間仍可能出現風格不一致或句子層級不自然的情形。
此外，單詞片段串接雖可形成可觀看動畫，但仍不等同於真人手語者在完整句子語境下的自然表達。
詞彙層級參數平滑雖能降低高頻抖動，但實驗亦顯示 hand-only PCK 下降，因此固定強度的平滑可能犧牲部分手形與快速手部動作細節。

第三，手語不僅依賴雙手動作，也包含臉部表情、視線、頭部姿態、身體方向與語氣等非手部訊息。
目前系統雖可保留部分追蹤資料中的臉部與上半身參數，但尚未針對疑問、否定、情緒或語氣等手語語法功能建立完整的非手部訊息生成機制。
因此，部分生成句子雖能呈現主要手部動作，但在自然手語語意與表情層面仍不夠完整。

第四，EHM-Tracker 追蹤結果會受到原始影片品質與 2D 偵測穩定度影響。
當手部模糊、遮擋、離開畫面或偵測信心較低時，追蹤結果可能產生手部偏移、抖動或誤追蹤。
此類問題會進一步影響後續 motion envelope 擷取、動作串接與 GUAVA 渲染結果。
因此，追蹤資料品質仍是整體系統效果的重要瓶頸。

最後，本研究之使用者問卷雖能反映系統在主觀觀看上的可接受程度，但填答者背景、手語能力與觀看經驗仍可能影響評分結果。
未來若要更完整評估系統於手語溝通或教學上的實用性，仍需納入更多手語使用者、手語教師或聽障族群，並設計更貼近實際使用情境的理解度與可用性測試。

\section{未來展望}

基於上述限制，未來研究可從語言、動作、追蹤資料與使用者互動等方向持續改善。

\begin{itemize}
    \item \textbf{擴充手語詞彙資料庫與語意對應能力}：
    未來可持續擴充已追蹤手語詞彙資料，並建立更完整的同義詞、複合詞與詞彙替代表。
    對於未收錄詞彙，亦可研究以語意相近詞彙、組合詞或拼字手語進行補足，使系統能處理更開放的中文輸入。

    \item \textbf{強化句子層級手語生成}：
    目前系統主要以詞彙片段串接形成句子動畫。
    未來可進一步加入句子層級的姿態規劃與速度調整，使多個 gloss 之間的動作不只是局部平滑，而能根據整句語意與節奏形成更自然的手語表達。

    \item \textbf{改善手部追蹤與手形穩定性}：
    使用者問卷與追蹤步數實驗皆顯示，手部細節是影響生成品質的重要因素。
    後續可針對低可信度手部偵測、手部遮擋與離畫面情境設計更穩定的篩選與修正策略，並加入手形一致性約束，以降低手指抖動與誤追蹤。參數平滑亦可依手部可見度、動作速度或不同身體部位採用自適應強度，在抑制抖動的同時保留快速手勢與手形細節。

    \item \textbf{加入非手部訊息生成機制}：
    手語中的疑問、否定、情緒與語氣常由臉部表情、頭部動作與身體姿態共同表達。
    未來可根據 SLT 模組輸出的語法標記或句型資訊，進一步生成對應的臉部表情與頭部姿態，使虛擬人動畫更接近自然手語。

    \item \textbf{建立自動化資料品質檢查流程}：
    本研究已使用 reprojection error、PCK 與視覺檢查分析追蹤品質。
    未來可將這些指標延伸為資料庫建立時的自動篩選機制，偵測手部離畫面、低信心關鍵點、手部抖動與異常追蹤片段，以提升大規模資料庫的穩定性。

    \item \textbf{提升系統互動性與輔助資訊呈現}：
    問卷回饋中有使用者建議加入文字或語音輔助。
    未來系統可於動畫播放時同步顯示中文句子、gloss sequence、詞彙片段來源或字幕資訊，使使用者更容易理解生成結果，也方便教學與錯誤分析。

    \item \textbf{擴大使用者評估與實際應用驗證}：
    後續可邀請更多手語使用者、聽障族群與手語教師參與評估，並加入理解度測驗、學習輔助效果與長期使用回饋，以更完整地驗證系統在手語學習、資訊近用與輔助溝通場景中的實用性。
\end{itemize}

總結而言，本研究完成了由手語動作資料建置至使用者動畫生成的模組化系統設計，並證明結合 SLT、既有手語影片資料、EHM-Tracker、動作參數串接與 GUAVA avatar rendering 之中文至虛擬人手語動畫流程具有可行性。
未來若能持續改善詞彙覆蓋率、追蹤資料品質、非手部訊息生成與句子層級自然度，將有助於推動更實用、更自然且更具可近性的虛擬人手語生成系統。

\end{document}
